\documentclass{article}
\usepackage[T1]{fontenc}
\usepackage{lmodern}
\usepackage{graphicx}
\usepackage{amsmath,amssymb}
\usepackage[a4paper,margin=2cm]{geometry}
\usepackage[absolute,overlay]{textpos}
\setlength{\TPHorizModule}{1mm}
\setlength{\TPVertModule}{1mm}
\usepackage[indonesian]{babel}
\usepackage{hyperref}
\setlength{\emergencystretch}{2em}
% unit: Halaman 16

\begin{document}

% === Halaman 16 (llm) ===

\section*{Enkripsi}

\begin{itemize}
    \item Enkripsi dilakukan dengan cara yang sama seperti algoritma \textit{knapsack} sebelumnya.
    \item Mula-mula plainteks dipecah menjadi blok bit yang panjangnya sama dengan kardinalitas barisan kunci publik.
    \item Kalikan setiap bit di dalam blok dengan elemen yang berkoresponden di dalam barisan kunci publik.
\end{itemize}

\end{document}
